\section{Built-in state constants}

State constants describe how a wire begins. They are not ordinary token names,
even though they participate in address resolution. Constants are
case-insensitive and may be written with a leading dollar sign. The three
supported base constants are \code{0p}, \code{1p}, and \code{sp}; the letter
\code{p} distinguishes protocol state notation from ordinary numbers.

\subsection{\code{0p}: the zero basis state}

\code{0p} denotes
\[
\ket0=\begin{bmatrix}1\\0\end{bmatrix}.
\]
\begin{lstlisting}
SET Output 0p
\end{lstlisting}
For workspace, this schedules zero preparation. The compiler batches pending
zero preparations into internal RESET operations immediately before the next
real gate, at an \code{INCREASECYCLE}, or at process end. RESET is deliberately
not source syntax; \code{SET Output 0p} is the public way to request zero.

Zero is the usual target initialization for reversible AND, OR, and XOR:
with $t=0$, the mutation $t\leftarrow t\oplus f(x)$ leaves $t=f(x)$.

\subsection{\code{1p}: the one basis state}

\code{1p} denotes
\[
\ket1=\begin{bmatrix}0\\1\end{bmatrix}.
\]
\begin{lstlisting}
SET Flag 1p
\end{lstlisting}
The compiler realizes one preparation by applying X to a zero-initialized
wire. A target beginning in one is useful when the desired reversible result is
the complement of a Boolean predicate. For example,
$1\oplus(A\land B)=\neg(A\land B)$, which is NAND.

\subsection{\code{sp}: equal superposition}

\code{sp} denotes the plus state
\[
\ket+=\frac{\ket0+\ket1}{\sqrt2}
=\frac1{\sqrt2}\begin{bmatrix}1\\1\end{bmatrix}.
\]
\begin{lstlisting}
SET Probe sp
\end{lstlisting}
The compiler realizes this preparation with a Hadamard gate. If immediately
measured, the result is 0 or 1 with equal probability. It is not merely a hidden
classical random bit: both amplitudes coexist and can interfere after later
phase and Hadamard operations.

In an expected truth-table \emph{output} cell, \code{sp} has a different
testing role: it means either measured bit is acceptable. Context therefore
matters---protocol SET initializes a plus state, while expected output metadata
uses \code{sp} as a don't-care marker.

\subsection{Dimension-suffixed constants}

The language also recognizes:
\begin{lstlisting}
0p_dim2
1p_dim8
sp_dim4
\end{lstlisting}
The suffix must be \code{_dim} followed by a power of two that is at least 2.
That number is a Hilbert-space dimension, so \code{_dim4} prepares two qubits
and \code{_dim8} prepares three. \code{0p} prepares every wire in $\ket0$,
\code{1p} flips only the least-significant wire, and \code{sp} applies H to
every wire. \code{_dim2} is the same one-wire preparation as the unsuffixed
constant. Any other dimension is rejected. Wire $k$ of register \code{R} is
named \code{R[k]}.

\subsection{Constants and state parameters}

For an ordinary local token, SET emits preparation behavior. For a PARAMS entry
typed \code{state}, SET supplies a default instead:
\begin{lstlisting}
PARAMS: Input:state
MAIN-PROCESS DefaultInput
SET Input 0p
X -I Input -O Input
RETURNVALS Input
\end{lstlisting}
The UI or truth-table runner may replace the default start state for Input.
This distinction is what allows one compiled protocol to be tested over every
binary input row instead of hard-wiring its parameter to zero.

